%% This is file `elsarticle-template-1-num.tex',
%%
%% Copyright 2009 Elsevier Ltd
%%
%% This file is part of the 'Elsarticle Bundle'.
%% ---------------------------------------------
%%
%% It may be distributed under the conditions of the LaTeX Project Public
%% License, either version 1.2 of this license or (at your option) any
%% later version.  The latest version of this license is in
%%    http://www.latex-project.org/lppl.txt
%% and version 1.2 or later is part of all distributions of LaTeX
%% version 1999/12/01 or later.
%%
%% Template article for Elsevier's document class `elsarticle'
%% with numbered style bibliographic references
%%
%% $Id: elsarticle-template-1-num.tex 149 2009-10-08 05:01:15Z rishi $
%% $URL: http://lenova.river-valley.com/svn/elsbst/trunk/elsarticle-template-1-num.tex $
%%
\documentclass[preprint,12pt]{elsarticle}
\makeatletter
\def\ps@pprintTitle{%
 \let\@oddhead\@empty
 \let\@evenhead\@empty
 \def\@oddfoot{}%
 \let\@evenfoot\@oddfoot}

%% Use the option review to obtain double line spacing
%% \documentclass[preprint,review,12pt]{elsarticle}

%% Use the options 1p,twocolumn; 3p; 3p,twocolumn; 5p; or 5p,twocolumn
%% for a journal layout:
%% \documentclass[final,1p,times]{elsarticle}
%% \documentclass[final,1p,times,twocolumn]{elsarticle}
%% \documentclass[final,3p,times]{elsarticle}
%% \documentclass[final,3p,times,twocolumn]{elsarticle}
%% \documentclass[final,5p,times]{elsarticle}
%% \documentclass[final,5p,times,twocolumn]{elsarticle}

%% The graphicx package provides the includegraphics command.
\usepackage{graphicx}
%% The amssymb package provides various useful mathematical symbols
\usepackage{amssymb}
%% The amsthm package provides extended theorem environments
%% \usepackage{amsthm}

%% The lineno packages adds line numbers. Start line numbering with
%% \begin{linenumbers}, end it with \end{linenumbers}. Or switch it on
%% for the whole article with \linenumbers after \end{frontmatter}.
\usepackage{lineno}

%% natbib.sty is loaded by default. However, natbib options can be
%% provided with \biboptions{...} command. Following options are
%% valid:

%%   round  -  round parentheses are used (default)
%%   square -  square brackets are used   [option]
%%   curly  -  curly braces are used      {option}
%%   angle  -  angle brackets are used    <option>
%%   semicolon  -  multiple citations separated by semi-colon
%%   colon  - same as semicolon, an earlier confusion
%%   comma  -  separated by comma
%%   numbers-  selects numerical citations
%%   super  -  numerical citations as superscripts
%%   sort   -  sorts multiple citations according to order in ref. list
%%   sort&compress   -  like sort, but also compresses numerical citations
%%   compress - compresses without sorting
%%
%% \biboptions{comma,round}

% \biboptions{}

\journal{Mike}

\begin{document}

\begin{frontmatter}

%% Title, authors and addresses

\title{Predicting Baseball Hall Of Fame Inductions}

%% use the tnoteref command within \title for footnotes;
%% use the tnotetext command for the associated footnote;
%% use the fnref command within \author or \address for footnotes;
%% use the fntext command for the associated footnote;
%% use the corref command within \author for corresponding author footnotes;
%% use the cortext command for the associated footnote;
%% use the ead command for the email address,
%% and the form \ead[url] for the home page:
%%
%% \title{Title\tnoteref{label1}}
%% \tnotetext[label1]{}
%% \author{Name\corref{cor1}\fnref{label2}}
%% \ead{email address}
%% \ead[url]{home page}
%% \fntext[label2]{}
%% \cortext[cor1]{}
%% \address{Address\fnref{label3}}
%% \fntext[label3]{}


%% use optional labels to link authors explicitly to addresses:
%% \author[label1,label2]{<author name>}
%% \address[label1]{<address>}
%% \address[label2]{<address>}

\author{Michael Hirsch}
\address{ILLC, University of Amsterdam}
\address{michaelahirsch@gmail.com}

\begin{abstract}
%% Text of abstract
Every year, the Baseball Writers Association of America votes on a new Hall of Fame class. Each ballot consists of 10 votes, and players need to appear on 75\% of ballots in order to be inducted into the hall. Players have 15 years to get inducted, and are no longer eligible if that time period has passed. Here we attempt to classify hall of fame batters based on their career statistics using decision trees and an ensemble method, random forests.
\end{abstract}


\end{frontmatter}

%%
%% Start line numbering here if you want
%%
%% main text
\section{Introduction}
\label{S:1}
Major League Baseball (MLB) has been keeping thorough records of batting, picthing, and fielding statistics since its innagural season in 1869. Recently, with the advent of sabermetrics by the Society for American Baseball Research (SABR), many  new metrics building on traditional statistics were created. Such metrics caught the eye of stasticians, as these new metrics allowed for the creation of even more powerful predictive models. 

In this paper, we are investigating what makes a Hall Of Fame (HOF) batter. There have been over 20,000 Major League baseball players in the history of the organization, and only 211 of them have been inducted into the HOF. There are several ways in which a player can be inducted, and we only concern ourselves with players inducted from Baseball Writers Association of America (BBWAA) ballots, as this process is regulated. Often, fans and players feel that a worthy candidate is unfairly denied entry into the HOF because of voter bias, stacked ballots, or negative associations with performance enhancing drugs. The model proposed in this paper should be able account for these cases.

\section{Data}
\label{S:2}
Data was collected from the Lahman Baseball Database, a freely available database that has been contiously updated since 1994 with the help of SABR and many individual researchers. There are many data tables available for download, but the ones that we are focusing on are \textit{Batting.dat} and \textit{HallOfFame.dat}. Batting statistics were supplied on a by-year basis, so a player's career statistics were computed by aggregating his yearly results. The batters' data was then merged with hall of fame voting results using the player's ID string. All data preparation was done within R, the environment where we also build our learning models. 

Attention is paid to 18 different batting statistics over the course of the player's career:

Below is a list of features used 
\begin{enumerate}
	\item G: games played
	\item AB: at-bats
	\item R: runs
	\item H: hits
	\item X2B: doubles
	\item X3B: triples
	\item HR: homeruns
	\item RBI: runs batted in
	\item SB: stolen bases
	\item CS: caught stealing
	\item BB: walks
	\item SO: strike outs
	\item IBB: intentional walks
	\item HBP: hit by pitches
	\item SH: sacrifice hits
	\item SF: sacrifice flys
	\item GIDP: grounded into double plays
	\item Inducted: classifier for HOF induction
\end{enumerate}

We needed to further filter our data by not including pitchers in our set of observations. This was done by removing those observations that indicated that the player recorded less than 300 RBI in his career. This number was not arbitrary, but rather chosen by cross referencing my data with HOF data from Baseball-Reference.


% here is where we can add citations
%Maecenas \cite{Smith:2012qr} fermentum \cite{Smith:2013jd} urna ac sapien tincidunt 


\section{Tree-Based Methods for Classification}

In this section, we will give an account of two learning methods that we will use for the basis of prediction: classification trees, and an ensemble method, random forests.

\subsection{Classification Trees}
Classification trees are a fairly effective predictive tool that are lauded for their high degree of interpretability. When constructing a classification tree, we begin with the full set of observations and begin dividing the predictor space into non-overlapping regions by means of binary splitting. In prediction, we assign each observation in a given region of the predictor space to the most commonly occuring class of the training observations in that region. When splitting, we are concerned with the textit{Information Gain} at each split, that is, we want to choose the feature to split on based on which feature split will give us the most pure sucessive node. This defined as:

$$IG(T,a) = H(T) - H(T|a)$$

\noindent where $T$ denotes a set of training examples, each of the form ($\textbf{x},y) = (x_{1},\ldots,x_{k},y)$ where each $x_a \in vals(a)$ is the value of the \textit{a}\textsuperscript{th} attribute of \textbf{x} and $y$ is the classification label. So, the information gain for a split on attribute \textit{a} is defined in terms of entropy as:

$$IG(T,a) = H(T) -\sum\limits_{v \in vals(a) } \frac{|\{\textbf{x} \in T : x_{a} = v\}|}{|T|} \cdot H(\{\textbf{x} \in T : x_{a} = v\})$$ 

\noindent Any decision tree algorithm will look too maximize this value, as that would indicate that the value of the attribute \textit{a} that we chose to split on will produce the most new information.

After splitting 





\subsection{Ensemble Methods}
\subsubsection{Bagging}
\subsubsection{Random Forests}

% \begin{table}[h]
% \centering
% \begin{tabular}{l l l}
% \hline
% \textbf{Treatments} & \textbf{Response 1} & \textbf{Response 2}\\
% \hline
% Treatment 1 & 0.0003262 & 0.562 \\
% Treatment 2 & 0.0015681 & 0.910 \\
% Treatment 3 & 0.0009271 & 0.296 \\
% \hline
% \end{tabular}
% \caption{Table caption}
% \end{table}

% \subsection{Subsection Two}

% Donec eget ligula venenatis est posuere eleifend in sit amet diam. Vestibulum sollicitudin mauris ac augue blandit ultricies. Nulla facilisi. Etiam ut turpis nunc. Praesent leo orci, tincidunt vitae feugiat eu, feugiat a massa. Duis mauris ipsum, tempor vel condimentum nec, suscipit non mi. Fusce quis urna dictum felis posuere sagittis ac sit amet erat. In in ultrices lectus. Nulla vitae ipsum lectus, a gravida erat. Etiam quam nisl, blandit ut porta in, accumsan a nibh. Phasellus sodales euismod dolor sit amet elementum. Phasellus varius placerat erat, nec gravida libero pellentesque id. Fusce nisi ante, euismod nec cursus at, suscipit a enim. Nulla facilisi.

% \begin{figure}[h]
% \centering\includegraphics[width=0.4\linewidth]{placeholder}
% \caption{Figure caption}
% \end{figure}

% Integer risus dui, condimentum et gravida vitae, adipiscing et enim. Aliquam erat volutpat. Pellentesque diam sapien, egestas eget gravida ut, tempor eu nulla. Vestibulum mollis pretium lacus eget venenatis. Fusce gravida nisl quis est molestie eu luctus ipsum pretium. Maecenas non eros lorem, vel adipiscing odio. Etiam dolor risus, mattis in pellentesque id, pellentesque eu nibh. Mauris nec ante at orci ultricies placerat ac non massa. Aenean imperdiet, ante eu sollicitudin vestibulum, dolor felis dapibus arcu, sit amet fermentum urna nibh sit amet mauris. Suspendisse adipiscing mollis dolor quis lobortis.

% \begin{equation}
% \label{eq:emc}
% e = mc^2
% \end{equation}

% \section{The Second Section}
% \label{S:2}

% Reference to Section \ref{S:1}. Etiam congue sollicitudin diam non porttitor. Etiam turpis nulla, auctor a pretium non, luctus quis ipsum. Fusce pretium gravida libero non accumsan. Donec eget augue ut nulla placerat hendrerit ac ut mi. Phasellus euismod ornare mollis. Proin tempus fringilla ultricies. Donec pretium feugiat libero quis convallis. Nam interdum ante sed magna congue eu semper tellus sagittis. Curabitur eu augue elit.

% Aenean eleifend purus et massa consequat facilisis. Etiam volutpat placerat dignissim. Ut nec nibh nulla. Aliquam erat volutpat. Nam at massa velit, eu malesuada augue. Maecenas sit amet nunc mauris. Maecenas eu ligula quis turpis molestie elementum nec at est. Sed adipiscing neque ac sapien viverra sit amet vestibulum arcu rhoncus.

% Vivamus pharetra nibh in orci euismod congue. Pellentesque habitant morbi tristique senectus et netus et malesuada fames ac turpis egestas. Quisque lacus diam, congue vel laoreet id, iaculis eu sapien. In id risus ac leo pellentesque pellentesque et in dui. Etiam tincidunt quam ut ante vestibulum ultricies. Nam at rutrum lectus. Aenean non justo tortor, nec mattis justo. Aliquam erat volutpat. Nullam ac viverra augue. In tempus venenatis nibh quis semper. Maecenas ac nisl eu ligula dictum lobortis. Sed lacus ante, tempor eu dictum eu, accumsan in velit. Integer accumsan convallis porttitor. Maecenas pretium tincidunt metus sit amet gravida. Maecenas pretium blandit felis, ac interdum ante semper sed.

% In auctor ultrices elit, vel feugiat ligula aliquam sed. Curabitur aliquam elit sed dui rhoncus consectetur. Cras elit ipsum, lobortis a tempor at, viverra vitae mi. Cras sed urna sed eros bibendum faucibus. Morbi vel leo orci, vel faucibus orci. Vivamus urna nisl, sodales vitae posuere in, tempus vel tellus. Donec magna est, luctus non commodo sit amet, placerat et enim.

%% The Appendices part is started with the command \appendix;
%% appendix sections are then done as normal sections
%% \appendix

%% \section{}
%% \label{}

%% References
%%
%% Following citation commands can be used in the body text:
%% Usage of \cite is as follows:
%%   \cite{key}          ==>>  [#]
%%   \cite[chap. 2]{key} ==>>  [#, chap. 2]
%%   \citet{key}         ==>>  Author [#]

%% References with bibTeX database:

\bibliographystyle{model1-num-names}
\bibliography{sample.bib}

%% Authors are advised to submit their bibtex database files. They are
%% requested to list a bibtex style file in the manuscript if they do
%% not want to use model1-num-names.bst.

%% References without bibTeX database:

% \begin{thebibliography}{00}

%% \bibitem must have the following form:
%%   \bibitem{key}...
%%

% \bibitem{}

% \end{thebibliography}


\end{document}

%%
%% End of file `elsarticle-template-1-num.tex'.